% !TeX root=../main.tex
% در این فایل، عنوان پایان‌نامه، مشخصات خود و چکیده پایان‌نامه را به انگلیسی، وارد کنید.

%%%%%%%%%%%%%%%%%%%%%%%%%%%%%%%%%%%%
\latinuniversity{University of Isfahan}
\latincollege{College of Computer Engineering}
\latinfaculty{Faculty of Computer Engineering}
\latindepartment{Artificial Intelligence and Robotics}
\latinsubject{Computer Engineering}
\latinfield{Artificial Intelligence and Robotics}
\latintitle{Developing a medical language model based on reasoning in Persian language}
\firstlatinsupervisor{Dr. Hamidreza Baradaran Kashani}
%\secondlatinsupervisor{Second Supervisor}
%\firstlatinadvisor{First Advisor}
%\secondlatinadvisor{Second Advisor}
\latinname{Mehrdad}
\latinsurname{Ghassabi}
\latinthesisdate{Feb 2026}
\latinkeywords{Artificial Intelligence in medicine,Persian language models,Medical language models,Natural language processing,Artificial Intelligence reasoning}
\en-abstract{
The use of artificial intelligence in answering medical questions has attracted significant attention in recent years as an emerging and important area in technology and healthcare. While extensive research in this field has been published in English, at the time this thesis began, the Persian language domain was lacking in research. In the first phase of this study, a medical corpus consisting of ninety million tokens and a dataset comprising twenty thousand medical questions and answers provided by physicians was collected from journals and online QA forums. Using part of this data, the base model "Aya-Expanse" was trained, and a new model named "Gaokerena-V" was introduced, which achieved 2.67\% better performance compared to its base model on the MMLU medical dataset translated into Persian; in the second phase, considering that medicine is heavily reliant on reasoning and logical analysis, a new model based on a chain of thoughts and logical reasoning was designed to significantly enhance its accuracy and efficiency. For this purpose, two new frameworks based on reinforcement learning with AI feedback were introduced to improve the reasoning capabilities of the base model, Aya-Expanse, through direct optimization of preferences. The result of this improvement was the "Gaokerena-R" model, which, despite being trained on a significantly smaller number of tokens compared to Gaokerena-V, achieved 6.34\% better performance than its base model on the same MMLU medical dataset translated into Persian. This level of improvement highlights the importance of reasoning in the field of medicine.
}
